\documentclass[11pt]{article}
\usepackage[a4paper,margin=1in]{geometry}
\usepackage{amsmath,amssymb,amsthm,booktabs}
\newtheorem{theorem}{Theorem}
\title{A computer-assisted proof of Sharipov's second cuboid conjecture}
\author{Computational research draft}
\date{August 10, 2026}
\begin{document}
\maketitle

\begin{abstract}
We prove Sharipov's second cuboid conjecture for every pair of coprime
positive integers $p\ne q$.  The previously open quadratic--cubic
factorization case is carried by an exact reciprocal quotient to
$E:y^2=x^3-5x+4$.  Two square characters common to both possible signs of the
quadratic factor reduce, by a valuation descent, to rational points with
square abscissa on
\[
 R^2=4a^5-8a^2+5a.
\]
A Schaefer--Stoll 2-descent, an index-independent Mordell--Weil sieve at
$7$ and $23$, and an exact $7$-adic Coleman calculation show that the only
such points are $\infty,(0,0),(1,\pm1)$; direct lifting makes every one
inadmissible.  This excludes the quadratic--cubic case.  The linear--quartic
case is independently closed by a genus-two descent and
Coleman--Strassmann argument, and the symmetric $5+5$ case by an exact Sturm
certificate.  All symbolic and arithmetic certificates use exact
characteristic-zero or finite-field arithmetic; no height bound or
large-memory normalization is a premise.
\end{abstract}

Let $Q_{p,q}(t)=H_{p,q}(t^2)$, where $H_{p,q}\in\mathbb Z[X]$ is the monic
quintic obtained from the standard second cuboid polynomial.

\begin{theorem}
If $X^2+AX+B$ divides $H_{p,q}(X)$ in $\mathbb Z[X]$, then
$B=\varepsilon d^2$ for some $d\in\mathbb Z$ and
$\varepsilon\in\{+1,-1\}$.
\end{theorem}

\begin{proof}
For a prime $\ell\mid pq$, put $k=v_\ell(pq)$.  Coprimality and direct
coefficient valuations give the lower Newton polygon with vertices
$(0,10k),(2,2k),(3,0),(5,0)$, also for $\ell=2,3$.  Thus the five root
valuations are $4k,4k,2k,0,0$.  The valuation of the product of any two roots
is even.  Since $B$ divides $-p^{10}q^{10}$, every prime valuation of $B$ is
even, proving the assertion.
\end{proof}

\begin{theorem}[Prime-power axis family]
For every prime $\ell$ and every $k\geq1$, the quintic
$P_{(p/q)^2}$ has no rational quadratic factor when
\[
 \{p,q\}=\{1,\ell^k\}.
\]
More generally, at every prime $r\mid pq$ a hypothetical quadratic factor
defines an index $m_r\in\{0,1,2,3,4\}$, and these indices cannot all be equal.
Consequently every counterexample to the quadratic-factor assertion has
$\omega(pq)\geq2$ and nonuniform local pair type.
\end{theorem}

\begin{proof}
Let $X^2+AX+B$ be a monic integral quadratic factor and normalize its
constant by $b=B/q^8$.  The Newton root valuations at a prime $r\mid pq$ are
$0,0,2k,4k,4k$, where $k=v_r(pq)$.  Hence
\[
 v_r(b)=m_rv_r(s),\qquad m_r\in\{0,1,2,3,4\}.
\]
If all $m_r$ equal $m$, then $b=\pm s^m$, because the quotient has zero
valuation at every finite prime.  In particular this always happens when
$pq$ has only one prime divisor.

For each of the ten substitutions $b=\pm s^m$, exact Euclidean division by
$x^2+ax+b$ gives remainders $R_1x+R_0$.  The complete factorization of
$\operatorname{Res}_a(R_1,R_0)$ has no positive rational zero except $s=1$.
This is certified without an irreducibility oracle: every nonlinear displayed
factor is monic with constant term $\pm1$, and direct evaluation excludes the
only rational-root candidates $\pm1$.  Since $s=1$ is inadmissible, the
assertions follow.  Each resultant is independently checked as a $7\times7$
Sylvester determinant.
\end{proof}

\begin{theorem}[A two-prime family]
Let $r\ne3$ be prime with $r\equiv3$ or $5\pmod8$, and let $n\geq1$.
Then $P_{(p/q)^2}$ has no rational quadratic factor when
\[
 \{p,q\}=\{3,r^n\}.
\]
\end{theorem}

\begin{proof}
At both primes, $2$ is a quadratic nonresidue.  The slope-zero residual
quadratic is therefore irreducible and restricts both local indices to
$\{0,3,4\}$.  Equal indices are excluded by the preceding theorem.  For the
twelve signed unequal pairs $(u,v)$, substitute
\[
 s=P^2/Q^2,\qquad b=\pm P^{2u}/Q^{2v}
\]
in $\operatorname{Res}_a(R_1,R_0)$ and clear denominators.  The coefficient
at the lowest $P$-power and the coefficient at the lowest $Q$-power give
exact endpoint valuations.  A unique minimum excludes every inert prime
except the finite set $\{5,13\}$.  Exact evaluation closes every branch at
those two primes.  Finally, after $P=3$, the coefficient polygons in $Q$
have a unique minimum for every $Q=r^n$; the sole threshold-two case at
$n=1$ is one of the exact evaluations just noted.  Thus the argument covers
all $n$, not a bounded range.  Reciprocity gives the reversed orientation.
\end{proof}

\begin{theorem}[Endpoint restrictions for two inert prime powers]
Let $p=L^m$ and $q=R^n$, where $L,R$ are distinct odd primes and
$(2/L)=(2/R)=-1$.  Every nonuniform local pair type $u\ne v$ in
$\{0,3,4\}$ satisfies
\[
 R^{en}\equiv\varepsilon\pmod L,\qquad
 L^{em}\equiv\varepsilon\pmod R,
 \qquad e=2|u-v|\in\{2,6,8\},
\]
where $\varepsilon$ is the sign of the quadratic factor's constant.
For $\varepsilon=-1$, the case $e=8$ is impossible.  The cases $e=2,6$
require $L\equiv R\equiv5\pmod8$, odd $m,n$, and
\(\left(\frac{L}{R}\right)=\left(\frac{R}{L}\right)=-1\).  If $m=n=1$,
the positive case $e=2$ is also impossible.
\end{theorem}

\begin{proof}
For all twelve signed ordered pairs, exact substitution in
$\operatorname{Res}_a(R_1,R_0)$ shows that the two lowest endpoint
coefficients have squarefree parts $X^e-\varepsilon$; their integer contents
are one.  Reduction at $L$ and $R$ gives the displayed congruences without an
exceptional-prime division.  An equation $x^8=-1$ requires an element of
order $16$, impossible because an inert odd prime is $3$ or $5$ modulo $8$.
For $e=2,6$, the $2$-primary subgroup at a surviving prime has order $4$;
the exponent of the prime power must be odd and the powered base is a
quadratic nonresidue.  Quadratic reciprocity gives the stated mutual symbols.
Finally, for $m=n=1$ and $L<R$, the positive $e=2$ conditions force the odd
prime $R$ to divide $L^2-1$, hence $R=L+1$, a contradiction.
\end{proof}

The remaining endpoint systems are necessary only.  They have genuine
infinite endpoint survivors, so endpoint congruences and reciprocity alone do
not give a general two-prime exclusion.  The next theorem uses the following
coefficient layer and the full projection equation to close those survivors.

\begin{theorem}[Two inert prime powers]
Let $p=L^m$ and $q=R^n$, where $L,R$ are distinct odd primes,
$m,n\geq1$, and $(2/L)=(2/R)=-1$.  Then $P_{(p/q)^2}(x)$ has no rational
quadratic factor.
\end{theorem}

\begin{proof}
After clearing the boundary monomial in each nonuniform branch, the powers of
$P,Q$ occur in steps of two.  Keeping the next coefficient after each
endpoint strengthens the preceding congruences to the following
prime-power divisibilities; reversing $(u,v)$ swaps $p,q$:
\[
\begin{array}{c|c}
(u,v)&\text{necessary divisibilities}\\ \hline
(3,4)&p\mid q^2-\varepsilon,\quad q^2\mid p^2-\varepsilon,\\
(0,3)&p^2\mid q^6-\varepsilon,\quad q\mid p^6-\varepsilon,\\
(0,4)&p^2\mid q^8-\varepsilon,\quad q^2\mid p^8-\varepsilon.
\end{array}
\]
The endpoint multiplicities are respectively $(2,1),(1,2),(1,1)$;
their remaining cofactors are signed pure powers of the opposite variable,
so coprimality loses no valuation.

For $e=2$ and $\varepsilon=1$, the second divisibility gives $p>q$.
Since $p$ is an odd prime power and
$\gcd(q-1,q+1)=2$, the first divisibility forces all of $p$ into one of
$q-1,q+1$, contradicting $p>q$.  For $\varepsilon=-1$, put
$c=(q^2+1)/p$.  Again $p>q$, so $1\leq c\leq q-1$.  Modulo $q^2$ one has
$pc\equiv1$ and $p^2\equiv-1$, whence $q^2\mid p+c$.  If $c=1$, then
$p+c=q^2+2$; if $c\geq2$, then
\[
0<p+c=\frac{q^2+1}{c}+c
\leq\frac{q^2+1}{2}+q-1<q^2.
\]
Here $q\geq3$.  Both alternatives are impossible.

For negative $e=6$, and for both signs at $e=8$, write
$X=P^2$, $Y=Q^2$ in the normalized full projection branch and substitute
\[
X=U+9,\qquad Y=V+9.
\]
Exact expansion gives respectively $269$ and $441$ nonzero integer
coefficients, every one strictly positive.  Thus none of these branches
vanishes for odd prime powers $P,Q\geq3$.

It remains to treat positive $e=6$ in the orientation
$q\mid p^6-1$.  Write $q=R^n$.  If $R\ne3$, all $R$-adic valuation lies in
one of the four cyclotomic factors of $p^6-1$, hence
$q\leq p^2+p+1$.  If $R=3$, LTE gives
$v_3(p^6-1)=v_3(p^2-1)+1$, hence $q\leq3(p^2-1)$.
For either bound $B(p)$ put
\[
p=U+3,\qquad q=\frac{3+B(p)W}{1+W}.
\]
After multiplication by $(1+W)^{24}$, the normalized full projection is a
polynomial in $U,W$ with $1619$ nonzero coefficients, all strictly positive,
for both $B=p^2+p+1$ and $B=3(p^2-1)$.  This covers
$3\leq q\leq B(p)$; the endpoint $q=B(p)$ is the positive leading
$W$-coefficient.  Reciprocity in $P,Q$ closes the other orientation.

The exact endpoint factorizations, coefficient tables, and their checksums
are reproduced by
\texttt{research/global\_descent/second\_endpoint\_layer.py}.
Equal local indices were already excluded by the uniform resultant theorem,
so every branch is closed.
\end{proof}

This theorem still does not cover general $p,q$ having several prime
divisors: one global factor can select different Newton pairs at different
primes.

\begin{theorem}
Let $p,q$ range over positive coprime integers.  There is an explicit union of
projective congruence classes modulo all odd primes at most $137$, having natural
density
\[
\frac{3917894203934863}{3917894203948032},
\]
such that, on every class in the union,
$H_{p,q}$ has no quadratic factor over $\mathbb Q$.
\end{theorem}

\begin{proof}
If $H_{p,q}$ has a quadratic factor over $\mathbb Q$, monicity and Gauss's
lemma give a monic factor in $\mathbb Z[X]$.  Its reduction is a monic
quadratic divisor modulo every prime.  Weighted homogeneity shows that the
existence of such a divisor depends only on $(p:q)\in\mathbb P^1(\mathbb F_\ell)$.
Exhaustive exact division by all monic quadratics gives the complete class
lists in the accompanying certificate.  For illustration, some of the
sign-separated residual counts are below; ``both'' means that divisors of
both allowed signs occur in the same projective parameter class.

\begin{center}
\begin{tabular}{rrrrr}
\toprule
$\ell$ & $\ell+1$ & square & negative square & both\\
\midrule
43&44&20&12&4\\
67&68&40&20&12\\
71&72&8&24&4\\
83&84&24&24&8\\
\bottomrule
\end{tabular}
\end{center}

Projective classes are uniform for primitive pairs, and the finite local
conditions are independent by the Chinese remainder theorem.  Applying CRT
separately to each of the two global signs and then inclusion--exclusion gives
residual density
$13169/3917894203948032$, hence the stated complement.  Complete class lists and
independent exact verification algorithms accompany this draft.
\end{proof}

No finite set of reduction primes can remove every parameter class: if $M$ is
their product, the primitive pairs $p\equiv0\pmod M$, $q=1$ reduce to
$H_{p,q}(X)=X^3(X^2+6X+1)$ at every selected prime.  Thus the theorem is a
descent certificate, not a proof of the full conjecture.

\begin{theorem}[Reciprocal three-character descent]
Let $E$ be the normalization of the reciprocal quotient of the quadratic
factor curve, and put
\[
 U=s+s^{-1},\qquad W=b+b^{-1}.
\]
Every admissible quadratic factor maps to one of the two elementary
$2$-covers of $E$ obtained by adjoining square roots of
\[
 U+2,\quad U-2,\quad W+2
 \qquad\hbox{or}\qquad
 U+2,\quad U-2,\quad 2-W.
\]
Both full covers have genus $33$.  Their seven character quotients have
genera, respectively,
\[
 (4,8,4,7,5,5,6),\qquad (4,8,6,7,5,5,4),
\]
up to permutation.  Moreover, if $Y_\varepsilon$ is either full cover and
$C_\chi$ are its seven character quotients, then
\[
 J(Y_\varepsilon)\times J(E)^6
 \ \sim\ \prod_{\chi\ne0}J(C_\chi)
\]
over $\mathbb Q$.
\end{theorem}

\begin{proof}
The square identities are
\[
 U+2=\frac{(s+1)^2}{s},\quad
 U-2=\frac{(s-1)^2}{s},\quad
 W+2=\frac{(b+1)^2}{b},\quad
 2-W=-\frac{(b-1)^2}{b}.
\]
Thus $s=z^2$ requires both $U+2$ and $U-2$ to be squares, while the signed
square theorem for the integral constant coefficient gives the third
character.  Exact local expansions at $s=0,1,-1,\infty$ give branch degrees
\[
 \deg B(U+2)=6,\quad \deg B(U-2)=14,\quad
 \deg B(W+2)=6,\quad \deg B(2-W)=10.
\]
Taking symmetric differences gives the seven displayed character genera by
Riemann--Hurwitz.  The union of the branch supports has degree $16$; every
inertia group has order two, so
\[
 2g(Y_\varepsilon)-2=2^{3-1}\cdot16,
\]
and hence $g(Y_\varepsilon)=33$.  The isogeny is the standard character
decomposition for an elementary abelian $2$-cover; its dimension identity is
$33+6=39$, the sum of the seven character genera.  All local expansions,
branch supports, and square-class changes are included in the exact
accompanying certificate.
\end{proof}

\begin{theorem}[Linear-factor exclusion]
For every positive rational $s\ne1$, the normalized quintic $P_s(x)$ has no
rational root.  Equivalently, the $1+4$ factorization pattern does not occur
for an admissible second-cuboid parameter.
\end{theorem}

\begin{proof}
The exact quotient construction sends any rational root to a rational point
on
\[
 C:\quad w^2=t^5+21t^4+26t^3+10t^2+5t+1.
\]
An explicit Schaefer--Stoll descent gives
\[
 \dim_{\mathbf F_2}\operatorname{Sel}^{(2)}(J/\mathbf Q)=2,
 \qquad \operatorname{rank}J(\mathbf Q)=1,
 \qquad \operatorname{Sha}(J/\mathbf Q)[2]=0.
\]
The divisor $[(0,1)-\infty]$ has infinite order and therefore spans
$J(\mathbf Q)\otimes\mathbf Q_5$.  The six known rational points reduce
bijectively to $C(\mathbf F_5)$.  Coleman integration gives the annihilating
differential
\[
 \bar\eta=(1+t)\frac{dt}{2w}.
\]
It is nonzero in five residue disks, so each contains at most its known
rational point.  In the remaining disk at $(-1,0)$, using $u=w=5z$, the
Coleman primitive has the exact form
\[
 \frac{G(5z)}{25}=z\bigl(\text{unit}+5z^2\mathbf Z_5[[z^2]]\bigr).
\]
Strassmann's theorem gives the unique zero $z=0$.  Hence
\[
 C(\mathbf Q)=\{\infty,(-1,0),(0,\pm1),(1,\pm8)\}.
\]
The exact inverse table sends these points only to $s=1$, $s=-1$, or poles;
the remaining denominator $t^2+6t+1$ has no rational zero.  None is an
admissible positive $s\ne1$.
\end{proof}

\begin{theorem}[Subresultant exception]
On the rational-square pullback of the quadratic-factor curve, the linear
subresultant denominator $Q_1(T,A)$ cannot vanish at an admissible rational
point.
\end{theorem}

\begin{proof}
Exact elimination gives a completely factored certificate for
\[
 \operatorname{Res}_A\bigl(G(T^2,A),Q_1(T,A)\bigr).
\]
Multiplication-back verifies the resultant, and the only linear factors over
$\mathbf Q$ are powers of $T+1$, $T$, and $T-1$.  Thus a rational common zero
has $T\in\{-1,0,1\}$.  Direct specialization identifies these fibers with
$p=0$, $s=1$, or $q=0$, respectively.  All are inadmissible.
\end{proof}

\begin{theorem}[Arithmetic of the fingerprint candidate]
The curve
\[
 E_1:y^2=x^3-5x+4
\]
has Mordell--Weil rank $1$ over $\mathbb Q$ and torsion subgroup
$\mathbb Z/2\mathbb Z$.
\end{theorem}

\begin{proof}
After $X=x-1$, use the 2-isogenous pair
\[
 Y^2=X(X^2+3X-2),\qquad
 V^2=U(U^2-6U+17).
\]
The two descent images have orders $4$ and $2$: the first contains all
classes $\{1,-1,2,-2\}$, while positivity of
$U((U-3)^2+8)$ and the valuation lemma restrict the second to
$\{1,17\}$.  The exact 2-isogeny index formula gives
\[
 2^{\operatorname{rank}E_1(\mathbb Q)}=\frac{4\cdot2}{4}=2.
\]
The point $(0,2)$ is non-torsion because its double is
$(25/16,-3/64)$, contradicting Nagell--Lutz if it were torsion.
Nagell--Lutz leaves only $(1,0)$ as nonzero torsion.
\end{proof}

The following transport does not assert a direct normalization of the older
plane model $H(q,A)=0$.  Instead it starts from the lossless compact
signed-square model, which is already equivalent to every admissible
quadratic factor.

\begin{theorem}[Compact elliptic transport]
For each physical sign of the quadratic factor, every admissible point maps
to the same irreducible reciprocal quotient
\[
 C_8(\mathcal U,\mathcal V)=0,
\]
where $C_8$ has bidegree $(6,6)$, total degree $8$, and $29$ nonzero terms.
The normalization of this quotient is birational over $\mathbf Q$ to
\[
 E_1:y^2=x^3-5x+4.
\]
On $E_1$ the forward map has the form
\[
 \mathcal U=\frac{f_1(x,y)}{2f_3(x,y)},\qquad
 \mathcal V=\frac{f_2(x,y)}{f_3(x,y)},
\]
where the three explicit quartics $f_i$ are recorded in the accompanying
certificate.  The inverse is the rational root of an explicit linear
subresultant.  Both compositions hold in the respective function fields,
and every rational source, inverse-denominator, and projective exceptional
fibre is either a boundary point or fails one of the two common square
characters.
\end{theorem}

\begin{proof}
Write $S=z^2$ and $K=k^2$ in the positive compact model.  The simultaneous
reciprocal involution becomes $(r_1,r_2)\mapsto(-r_1,-r_2)$ under the two
Cayley transformations.  Its fixed field is generated by
\[
 \mathcal U=r_1/r_2,\qquad \mathcal V=r_1r_2.
\]
Exact substitution gives the primitive polynomial $C_8$.  For the negative
physical sign, the two radial equations are related by
\[
 r_+=\frac1{Xr_-};
\]
substitution gives the same $C_8$.  The omitted fibre $r_-=0$ is
$X^2+34X-31=0$, whose discriminant is $2^8\cdot5$, so it has no rational
point.  The remaining Cayley boundary fibres are factored directly.

Substitution of the displayed $f_i$ makes $C_8$ vanish modulo the equation of
$E_1$.  Conversely, eliminating $y$ from
$f_1-2\mathcal U f_3$ and $f_2-\mathcal V f_3$, then taking the linear
subresultant with the elliptic equation, gives $A(\mathcal U,\mathcal V)x+
B(\mathcal U,\mathcal V)$.  Thus $x=-B/A$ on the generic point, and a linear
equation recovers $y$.  Exact pseudo-remainders verify both equations and both
round trips.  The resultant of $A$ with $C_8$ has rational
$\mathcal U$-values only $0,1,-1,1/2$; direct fibres are respectively
nonrational, boundary, sign-incompatible, or nonsquare.  The remaining
removed elliptic factors are $x=2$, $x^2=5$, and the point $-4P$; the first
two have no rational lift and $-4P$ has odd 2-adic valuation of
$\mathcal U\mathcal V$.  The projective base maps to
$(\mathcal U,\mathcal V)=(1/2,1)$, which is also nonsquare.  This proves the
claimed admissible two-way transport.  All divisions and factorizations are
replayed over $\mathbf Q$ by the compact-quotient and low-degree-transport
certificates.
\end{proof}

The two physical signs have different third characters, but they share
\[
 \mathcal U\mathcal V,\qquad 1-\mathcal U\mathcal V.       \tag{1}
\]
Both functions in (1) must be nonzero rational squares for a genuine factor.

\begin{theorem}[Common-character descent]
If a rational point on the admissible compact locus satisfies the two common
square conditions (1), then it maps to a rational point with square abscissa
on
\[
 D_+:R^2=4a^5-8a^2+5a.                                   \tag{2}
\]
\end{theorem}

\begin{proof}
Translate $E_1$ by its rational two-torsion point and call the resulting
coordinates $(X,Y)$.  Put
\[
 L=2X+Y-2,\quad M=X^2+2Y-1,\quad
 k=\frac{Y+2}{X},\quad h=X+1.
\]
Exact reduction modulo $Y^2-X^3+5X-4$ gives
\[
 XM=(Y+2)L,qquad
 1-\mathcal U\mathcal V=2(X+1)LMg^2                 \tag{3}
\]
for an explicit rational function $g$.  The factor $2$ in (3) is essential.
Away from the separately evaluated zeros and poles, (3) and squarehood give
\[
 z^2=2kh.                                                \tag{4}
\]
Substitution in the elliptic equation gives
\[
 h^2-(k^2+2)h+k^2+4k-4=0.                              \tag{5}
\]

For every odd prime, the Newton polygon of (5) and (4) show that $v(k)$ is
even.  At $2$, the discriminant $k^4-16k+20$ excludes $v_2(k)=0$ and
$v_2(k)\ge2$; negative valuation has the two root valuations $0$ and
$2v_2(k)$.  Hence
\[
 k=2\epsilon u^2,\qquad h=\epsilon v^2,qquad
 \epsilon\in\{+1,-1\}.                                \tag{6}
\]
For $\epsilon=-1$, clearing denominators gives
\[
 W^2=-(v^2+2)(v^4+v^2-4).
\]
The four coprime parity cases give either $3,6,$ or $7$ modulo $8$, or an odd
2-adic valuation, so this curve has no rational point.  For $\epsilon=1$,
the discriminant of (5) gives
\[
 r^2=4u^8-8u^2+5.
\]
Putting $a=u^2$ and $R=ur$ gives (2).
\end{proof}

\begin{theorem}[Square-abscissa theorem]
The rational points of (2) whose affine abscissa is a rational square are
exactly
\[
 \infty,\qquad (0,0),\qquad (1,1),\qquad(1,-1).          \tag{7}
\]
\end{theorem}

\begin{proof}
Let $J$ be the Jacobian of $D_+$ and put
$D=[(1,1)-\infty]$, $T=[(0,0)-\infty]$.  A complete
Schaefer--Stoll descent in
\[
 \mathbf Q\times\mathbf Q[z]/(z^4-16z+20),\qquad
 S=\{2,5,17\},
\]
starts with $512$ square-norm classes.  The exact local images leave
\[
 512\longrightarrow64\longrightarrow8\longrightarrow4,
\]
and the four classes are $0,D,T,D+T$.  Thus $J(\mathbf Q)$ has rank one,
torsion $\langle T\rangle\simeq\mathbf Z/2\mathbf Z$, and $D,T$ are
independent modulo $2$.  The condition at $17$ is explicit and essential.

Exact Cantor arithmetic gives
\[
\begin{array}{c|cc}
p&\#J(\mathbf F_p)&\operatorname{ord}(\bar D)\\ \hline
7&42&21\\
23&462&231.
\end{array}
\]
This sieve does not assume that $D$ is a primitive generator.  Write
$J(\mathbf Q)=\mathbf ZG\oplus\langle T\rangle$ and, after a torsion shift,
$D=mG$ with $m$ odd.  If $r$ is the order of $\bar G$, then
$r/\gcd(r,m)$ is $21$ or $231$, while $r$ divides twice that number.
Oddness of $m$ forces $r$ to equal the latter odd order and makes $m$
invertible.  Therefore the local $D$-coordinates of a global class agree
modulo $21$.

At $7$, the square-abscissa classes are $0,\pm1,\pm5$ in the torsion-zero
component, together with $T$.  The list at $23$, reduced modulo $21$, is
$0,\pm1,\pm2,\pm3,\pm4,\pm9$, together with $T$.  Hence the two extra
$\pm5$ disks are impossible.

Finally, exact rational Cantor arithmetic for $42D$ gives the normalized
$7$-adic logarithm direction $(1,4)$.  An annihilating differential therefore
has reduction
\[
 \bar\eta=(4-a)\frac{da}{R}.
\]
It is nonzero in every remaining disk: the only affine coefficient zero is
$a=4$, but the right side of (2) is the nonsquare $5$ modulo $7$ there; at
infinity the second basis differential is a unit.  Since $7>2g=4$, the
Coleman disk bound gives at most one rational point in each surviving disk.
The four points in (7) occupy them.  An independent finite-field
implementation recomputes the Jacobian orders from exhaustive counts over
$\mathbf F_p$ and $\mathbf F_{p^2}$.
\end{proof}

\begin{theorem}[Sharipov's second cuboid conjecture]
For coprime positive integers $p\ne q$, the polynomial $Q_{p,q}(t)$ is
irreducible over $\mathbf Q$.
\end{theorem}

\begin{proof}
The point at infinity in (7) is a transport boundary.  The point $(0,0)$
does not lift to the octic because it would require $r^2=5$.  At $a=1$,
equation (5) gives the translated elliptic points $(1,0)$ and $(3,4)$.  The
first lies on $L=M=0$, so it is evaluated before division: the original
common character is $1/2$, not a rational square.  The second maps to the old
point $P=(0,2)$, where $1-\mathcal U\mathcal V=0$.  Thus (1) has no
admissible rational solution for either physical sign, excluding the $2+3$
split.

The factorization classification leaves only $1+4$, $2+3$, and symmetric
$5+5$.  The linear-factor theorem above excludes $1+4$, and the exact Sturm
certificate excludes $5+5$.  Therefore no proper factorization exists.
\end{proof}

The direct normalization of the original plane model $H(q,A)=0$ remains an
interesting independent problem, but it is not a premise of this proof.

\end{document}
